% Assume-se que \pretextual já foi feito

\imprimircapa%
\imprimirfolhaderosto*
% \protect\incluirfichacatalografica{ficha.pdf}  % descomentar quando tiver a ficha
{}% placeholder: evita que \imprimirfolhadecertificacao seja engolido pelo \@ifstar
\imprimirfolhadecertificacao


% \begin{dedicatoria}
% \end{dedicatoria}


\begin{agradecimentos}
Ao Prof.\ Márcio Bastos Castro, pela orientação cuidadosa, pela liberdade de
exploração e pelo ambiente de pesquisa que tornou este trabalho possível.

Ao LaPeSD, em especial a Vanderlei Munhoz, cujo desenvolvimento do orquestrador
HPC@Cloud serviu de base para a integração proposta neste trabalho, e cujas
contribuições técnicas ao longo do projeto foram determinantes.

Este trabalho foi parcialmente financiado pelo Conselho Nacional de
Desenvolvimento Científico e Tecnológico (CNPq) e pela Amazon Web Services
(AWS) por meio da Chamada CNPq/AWS Nº~64/2022 (Cloud Credits for Research).

À minha família, pelo apoio incondicional ao longo de toda a graduação.
\end{agradecimentos}


% \begin{epigrafe}
% \end{epigrafe}


\begin{resumo}[Resumo]
Instâncias \textit{spot} da Amazon Web Services oferecem descontos de até 90\%
em relação ao preço convencional, mas podem ser revogadas pelo provedor a
qualquer momento, com aviso de apenas dois minutos, descartando todo o progresso
computacional acumulado em aplicações MPI de longa duração. Este trabalho propõe
e avalia uma arquitetura de execução resiliente para aplicações MPI legadas em
instâncias \textit{spot}, sem exigir modificação de código, por meio da
integração do MANA, sistema de \textit{checkpoint/restart} transparente para MPI,
ao orquestrador HPC@Cloud. A topologia adotada é híbrida: nó coordenador em
instância \textit{burstable on-demand} e nós de computação em instâncias
\textit{spot}. Duas estratégias de recuperação são implementadas: a
\textsc{Replace}, que reprovísiona o nó interrompido e retoma a execução com a
capacidade original do \textit{cluster}, e a \textsc{Degraded}, que marca o nó
como inativo e reinicia imediatamente com $N{-}1$ \textit{workers}. A avaliação
experimental, conduzida com 282 execuções sobre os \textit{benchmarks} NPB
CG-C, EP-D e LU-C em \textit{clusters} de 2, 4 e 8 \textit{workers}, revelou
dois regimes de viabilidade: para \textit{jobs} curtos como o CG-C (12 a 35~s),
o \textit{overhead} de recuperação domina e a abordagem não se justifica; para
\textit{jobs} de duração moderada a longa como o EP-D e o LU-C, a
\textsc{Degraded} alcançou menor tempo de execução em 17 das 18 configurações
avaliadas, com economia de custo de 2 a 38\% em relação à execução equivalente
em instâncias \textit{on-demand} sem tolerância a falhas.

\vspace{\baselineskip}
\textbf{Palavras-chave:} Computação em Nuvem. Computação de Alto Desempenho.
Instâncias \textit{Spot}. Tolerância a Falhas. \textit{Checkpoint/Restart}.
MANA. HPC@Cloud.
\end{resumo}


\begin{resumo}[Resumo Estendido]

\section*{Introdução}

A computação de alto desempenho (HPC) tem migrado progressivamente para
ambientes de nuvem pública, atraída pela elasticidade de recursos e pelo modelo
de cobrança por uso. Nesse contexto, as instâncias \textit{spot} da Amazon Web
Services representam uma oportunidade econômica relevante: ao adquirir
capacidade computacional ociosa do provedor, é possível obter descontos de até
90\% em relação ao preço convencional. Essa vantagem de custo, contudo, vem
acompanhada de uma restrição operacional significativa: o provedor pode revogar
as instâncias \textit{spot} a qualquer momento, com aviso mínimo de dois
minutos antes da terminação efetiva.

Para aplicações de curta duração, a probabilidade de interrupção é baixa e o
impacto de uma revogação é tolerável. Para aplicações MPI de longa duração,
porém, uma revogação descarta todo o progresso computacional acumulado, forçando
a reexecução completa do \textit{job}. Essa característica torna as instâncias
\textit{spot} impraticáveis para esse tipo de carga sem um mecanismo de
preservação de estado, especialmente quando as aplicações são legadas e não
preveem modificação de código.

O MANA (\textit{MPI-Agnostic Network-Agnostic checkpoint/restart}) é um sistema
de \textit{checkpoint} transparente para aplicações MPI que opera sem qualquer
modificação no código da aplicação. Baseado no DMTCP (\textit{Distributed
MultiThreaded CheckPointing}), o MANA captura a imagem completa do estado de
execução de todos os processos MPI e permite retomá-la a partir do ponto salvo,
mesmo após uma falha de nó. Essa propriedade cria a oportunidade de combinar a
economia das instâncias \textit{spot} com a resiliência necessária para
\textit{jobs} de longa duração.

Este trabalho integra o MANA ao orquestrador HPC@Cloud para construir uma
arquitetura de execução resiliente e econômica para aplicações MPI legadas em
instâncias \textit{spot}. A questão central investigada é se e quando essa
combinação é viável, tanto em termos de tempo de execução quanto de custo
monetário, em comparação com a alternativa de instâncias \textit{on-demand} sem
tolerância a falhas.


\section*{Objetivos}

O objetivo geral deste trabalho é propor, implementar e avaliar uma arquitetura
de execução resiliente e econômica para aplicações HPC em instâncias
\textit{spot} da AWS, integrando o MANA ao HPC@Cloud e quantificando o impacto
em desempenho e custo em comparação com execução nativa sem tolerância a falhas.

Os objetivos específicos são:
\begin{enumerate}
  \item Consolidar a análise de viabilidade de instâncias \textit{burstable}
        para HPC em nuvem, identificando seus limites operacionais e justificando
        seu uso restrito ao papel de coordenador de \textit{checkpoint};
  \item Estender o HPC@Cloud com suporte à topologia híbrida
        \textit{head}/\textit{worker}, incluindo diferenciação de papéis por
        tipo de instância e as implicações no processo de provisionamento;
  \item Implementar o \textit{bootstrap} automatizado do ambiente MANA/DMTCP
        sobre \textit{clusters} AWS, incluindo configuração do Slurm,
        implantação dos binários MANA e sincronização de estado via EFS;
  \item Implementar um mecanismo de detecção de interrupção e duas estratégias
        de recuperação (\textsc{Replace} e \textsc{Degraded}), com suporte a
        falhas sintéticas reproduzíveis para fins de avaliação controlada;
  \item Quantificar o \textit{overhead} do MANA sobre os NAS Parallel
        Benchmarks (NPB), especificamente CG-C, EP-D e LU-C, em
        \textit{clusters} com 2, 4 e 8 \textit{workers}, incluindo estudos
        sintéticos para isolar os mecanismos de \textit{overhead};
  \item Comparar o custo-benefício das estratégias \textsc{Replace} e
        \textsc{Degraded} em relação à execução nativa, identificando as
        condições sob as quais cada estratégia é economicamente vantajosa.
\end{enumerate}


\section*{Metodologia}

A topologia adotada é híbrida: um nó coordenador (\textit{head}) opera como
instância t3.large \textit{on-demand}, imune a revogações, e os nós de
computação (\textit{workers}) executam como instâncias m5.xlarge \textit{spot},
aproveitando os descontos de até 90\%. Essa escolha foi fundamentada na análise
de viabilidade de instâncias \textit{burstable} conduzida previamente, que
demonstrou que o perfil intermitente do coordenador de \textit{checkpoint} é
compatível com as limitações de CPU e rede dessas instâncias, enquanto cargas
HPC com sincronização frequente sofrem degradação expressiva.

A integração ao HPC@Cloud implementa três componentes. O \textit{bootstrap}
automatizado configura o ambiente Slurm e implanta os binários do MANA no EFS
compartilhado, tornando os nós operacionais sem intervenção manual. O
\textit{watcher} monitora periodicamente o estado das instâncias \textit{spot}
via API da EC2 e, ao detectar uma revogação iminente, aciona o salvamento do
estado de todos os processos MPI. Após a confirmação do
\textit{checkpoint}, o sistema aplica uma das duas estratégias configuráveis: a
\textsc{Replace}, que reprovísiona o nó interrompido e retoma o \textit{job}
com a capacidade original do \textit{cluster}, e a \textsc{Degraded}, que marca
o nó como inativo e reinicia imediatamente com $N{-}1$ \textit{workers}.

Para isolar cada componente do custo de recuperação, a execução é decomposta em
quatro fases: P0 (computação antes da falha), P1 (escrita do
\textit{checkpoint}), P2 (reconfiguração do nó, subdividida em P2a$+$P2c
comuns às duas estratégias e P2b onde elas divergem) e P3 (computação após o
reinício). Essa decomposição permite atribuir cada segundo observado a uma causa
específica, separando tempo de computação de tempo de infraestrutura e de
coordenação.

A avaliação experimental utilizou os \textit{benchmarks} NPB CG-C
(\textit{Conjugate Gradient}, classe~C), EP-D (\textit{Embarrassingly Parallel},
classe~D) e LU-C (\textit{Lower-Upper Gauss-Seidel}, classe~C), representando
respectivamente regimes de comunicação intensiva com buffers por processo,
operações coletivas sem estado por \textit{rank} e comunicação estruturada com
vizinhos de cardinalidade fixa. A matriz de experimentos combinou 3
\textit{benchmarks} por 3 tamanhos de \textit{cluster} (2, 4 e 8
\textit{workers}) por 3 momentos de falha (10\%, 25\% e 50\% do tempo total)
por 2 estratégias, além das linhas de base noFT e MANA-noFT, totalizando 282
execuções válidas com $N{=}3$ repetições por configuração. A injeção de falhas
foi realizada de forma sintética e reproduzível por meio do mecanismo
\texttt{auto\_test\_failure} do HPC@Cloud.


\section*{Resultados e Discussão}

A avaliação revelou que o \textit{overhead} do MANA em condições normais
(sem falhas) é aplicação-dependente e emerge da interação entre o padrão de
execução de cada aplicação e a camada de instrumentação do MANA: varia de 27\%
no EP-D com 8~\textit{workers} a 222\% no CG-C com 2~\textit{workers}. Os
estudos sintéticos identificaram que a frequência de chamadas MPI não é o
mecanismo amplificador; o custo intrínseco medido em programas isolados é de
apenas 4 a 5~s por execução. A distância entre esse valor e os percentuais
observados nos \textit{benchmarks} reais indica que o \textit{overhead} relevante
é específico ao padrão de execução de cada aplicação.

A análise das fases de recuperação revelou onde as estratégias divergem: a fase
P2b, que compreende a reconfiguração do nó após a falha. A \textsc{Degraded}
conclui P2b em $\approx$11~s ao atualizar o estado do nó no Slurm; a
\textsc{Replace} aguarda $\approx$120~s pelo reaprovisionamento de um nó
substituto na AWS. Essa diferença de $\approx$109~s é o elemento central da
comparação entre estratégias, e seu peso relativo é modulado por dois fatores
que interagem: o tamanho do \textit{cluster} ($N$) e o volume de trabalho
remanescente no momento da falha.

Em \textit{clusters} maiores, a perda de um nó representa uma fração menor da
capacidade total, reduzindo a penalidade de P3 da \textsc{Degraded} e tornando
a economia de P2b dominante. Em \textit{clusters} pequenos com falha precoce, o
custo de P3 com $N{-}1$ \textit{workers} pode superar os $\approx$109~s
economizados, tornando a \textsc{Replace} mais eficiente. Essa estrutura de
\textit{crossover} foi observada empiricamente em apenas um dos 18 cenários
avaliados para EP-D e LU-C: o EP-D com 2~\textit{workers} e falha a 10\% do
tempo total. Nas demais 17 configurações, a \textsc{Degraded} apresentou menor
tempo de execução, com vantagem crescente conforme $N$ aumenta.

Os resultados revelam dois regimes econômicos distintos. Para \textit{jobs}
curtos como o CG-C (12 a 35~s), o \textit{overhead} de recuperação equivale a
3 a 17 vezes o tempo útil de cálculo, e o desconto \textit{spot} não compensa
nem em tempo nem em custo. Para \textit{jobs} de duração moderada a longa como
o EP-D e o LU-C, a \textsc{Degraded} alcança economia de 2 a 38\% em relação
ao equivalente \textit{on-demand} sem recuperação, com as análises de tempo e
de custo apontando consistentemente na mesma direção. A viabilidade depende,
portanto, de forma determinante da duração do \textit{job}: apenas quando o
custo fixo de recuperação representa uma fração minoritária do tempo total é que
a combinação de \textit{spot} com tolerância a falhas se torna vantajosa.


\section*{Considerações Finais}

A integração proposta demonstra que execução resiliente e econômica de
aplicações MPI legadas em instâncias \textit{spot} é alcançável sem modificação
de código, desde que a duração do \textit{job} seja suficiente para amortizar o
custo de recuperação. Os seis objetivos específicos foram atendidos: a análise
de instâncias \textit{burstable} fundamentou a topologia adotada; o HPC@Cloud
foi estendido com \textit{bootstrap} automatizado e dois mecanismos de
recuperação configuráveis; e a avaliação experimental quantificou o
\textit{overhead}, caracterizou o comportamento das estratégias e identificou as
condições de viabilidade econômica.

Como extensões naturais do trabalho, destacam-se: a validação empírica do regime
de \textit{jobs} de longa duração, com classes mais custosas dos
\textit{benchmarks}, observando a restrição do limite de dois minutos para
escrita do \textit{checkpoint}; a avaliação do comportamento das estratégias
sob múltiplas falhas sequenciais; e a adição de \textit{checkpoints} periódicos
combinados com predição de revogações baseada em histórico de interrupções de
instâncias \textit{spot}.

\vspace{\baselineskip}
\textbf{Palavras-chave:} Computação em Nuvem. Computação de Alto Desempenho.
Instâncias \textit{Spot}. Tolerância a Falhas. \textit{Checkpoint/Restart}.
MANA. HPC@Cloud.
\end{resumo}


\begin{abstract}
AWS spot instances offer discounts of up to 90\% compared to on-demand pricing,
but may be revoked by the provider at any time with only two minutes of notice,
discarding all accumulated computational progress in long-running MPI
applications. This work proposes and evaluates a resilient execution
architecture for legacy MPI applications on AWS spot instances, without
requiring code modification, by integrating MANA, a transparent
\textit{checkpoint/restart} system for MPI, with the HPC@Cloud orchestrator.
The adopted topology is hybrid: a coordinator node on a \textit{burstable
on-demand} instance and compute nodes on \textit{spot} instances. Two recovery
strategies are implemented: \textsc{Replace}, which reprovisions the interrupted
node and resumes execution with the original cluster capacity, and
\textsc{Degraded}, which marks the node as inactive and immediately restarts
with $N{-}1$ workers. The experimental evaluation, conducted with 282 executions
on the NPB benchmarks CG-C, EP-D, and LU-C on clusters of 2, 4, and 8 workers,
revealed two viability regimes: for short jobs such as CG-C (12 to 35~s),
recovery overhead dominates and the approach is not worthwhile; for medium to
long-running jobs such as EP-D and LU-C, \textsc{Degraded} achieved shorter
execution times in 17 out of 18 evaluated configurations, with cost savings of
2 to 38\% compared to the equivalent on-demand execution without fault
tolerance.

\vspace{\baselineskip}
\textbf{Keywords:} Cloud Computing. High Performance Computing. Spot Instances.
Fault Tolerance. Checkpoint/Restart. MANA. HPC@Cloud.
\end{abstract}


\listoffigures*
\listoftables*
\listofalgorithms*

%%% Força todas as siglas definidas em acronyms.tex a aparecer na lista,
%%% mesmo sem uso de \gls{} no texto.
\glsaddall
\listadesiglas*[7em]

\tableofcontents*%

%%% Local Variables:
%%% mode: latex
%%% TeX-master: "main"
%%% End:
